\chapter{Technical Implementation}\label{ch:method}

This chapter outlines the technical implementation of the proposed concept, with a particular focus on the analysis and modification of the 3MF file structure and its integration into the \textit{OrcaSlicer} environment. 
The goal was to enable the transfer of model appearances from Fusion to OrcaSlicer in a manner that preserves and utilises colour-based metadata as print hints.
To demonstrate the feasibility of the concept, the implementation focused on enabling the application of presets to components based on their associated colours during import.
For this purpose, \textit{OrcaSlicer}'s code repository was forked, using Git as a version control system. 
The subsequent sections illustrate the process of analysing \textit{Fusion}’s 3MF export behaviour, identifying relevant colour data structures, adapting \textit{OrcaSlicer}'s import logic, and extending the user interface to support the proposed functionality.

\section{Technical Analysis of the 3MF File Structure}

Before the implementation, a comprehensive analysis of the 3MF file format was conducted to determine how \textit{Fusion} stores appearance data. 
The 3MF format organises the main model-related data in hierarchical XML structures.

Per the additional Materials and Properties Extension specification\cite{3mf_consortium_materials_2025}, colour information can be encoded in several different ways.
One such option is the colour group, which \textit{Fusion} utilises. 
Each group is identified by a unique parameter ID, referenced per object as \texttt{pid}, and contains one or more \texttt{<color>} entries
The colours are defined using an RGBA colour code in hexadecimal format and are implicitly assigned a 0-based index. 
In compliance with the specification, the alpha channel is encoded but not utilised, and can therefore be disregarded in this context.
In the file structure example \ref{lst:3dmodel.model}, these sections can be found in lines 4-40 (\texttt{<resources>}), enclosing the \texttt{<m:colorgroup>} and \texttt{<object>} tags.

\subsection{\textit{Fusion}'s Export Behaviour}

An initial uncertainty concerned whether \textit{Fusion} exported appearance data by default, as no colours appeared during import into \textit{OrcaSlicer}. 
Upon closer inspection, it became evident that Fusion does, in fact, include appearance information; however, each appearance is represented as a colour group containing two colour entries. 
The first entry typically approximates the visual appearance assigned in \textit{Fusion}, while the second entry is consistently a light grey tone.

Assigning multiple appearances to a single component merely resulted in additional \texttt{<color>} tags within the same group, without altering the pattern that the final colour remained grey. 
Since Fusion allows appearances and physical materials to be applied independently, it was initially hypothesised that one colour might represent the appearance while the other corresponded to the physical material. 
However, neither documentation nor experimental evidence substantiates this assumption, and the precise purpose of the second colour remains undetermined.


\section{Analysis \textit{OrcaSlicer}'s 3MF Import Logic}

In order to understand why \textit{Fusion}’s exported colours were not displayed, the source code of \textit{OrcaSlicer} was examined thoroughly. 
Preliminary analysis indicated that the function responsible for parsing XML colour data was not being executed. 
However, upon closer inspection it became apparent that the function was indeed called and populated data structures — albeit not the ones represented in the Model class (\ref{fig:model_class_uml}).

\begin{figure}
    % \vspace{-1pt}
  \includegraphics[width=\linewidth]{Model.png}
  \label{fig:model_class_uml}
  \caption{Excerpt of \textit{OrcaSlicer}'s Model class with Relevant Attributes highlighted}
\end{figure} 

The observed behaviour could be traced to the handling of colour groups during import. 
The colours within each group were processed sequentially, but only the last colour in the group was retained in memory, effectively overwriting the others. 
This behaviour explains why only a single (typically grey) colour appeared after import.

To address this, the import logic was modified to preserve the \textit{first} colour in each group instead of the \textit{last}. 
Following this adjustment, all colours defined in the 3MF file were correctly interpreted and handled within OrcaSlicer. 
This modification enabled the application to assign one filament preset to each colour detected across multiple model components.
Specifically, it applied the user's saved filament presets in the saved order consecutively to each color it encountered. 

\subsection{Internal Data Handling}

After confirming that colours could be successfully imported, the next step was to identify the location of this information within \textit{OrcaSlicer}’s internal model representation. 
Contrary to initial assumptions, the colour attributes were not associated with the material attributes of individual model instances (\texttt{ModelMaterialMap}). 

Instead, the data was found to be associated with the extruder configurations, which handle the assignment of filaments and their corresponding parameters.
In \textit{OrcaSlicer}, extruder information is implemented as a general ConfigOption object (see \ref{fig:model_class_uml}), a wrapper class encapsulating a key-value pair with support for multiple variable types.
Since this structure abstracts the data type without explicitly defining its contents, identifying and tracing the relevant information introduced additional complexity to the analysis of the underlying codebase.

\section{User Interface Adaption}

The user interface design for importing 3MF models with colour mappings was inspired by the OBJ import workflow, which allows users to map model colours to existing filament presets, as shown in \ref{fig:obj_color_dialog}. 
Therefore, the corresponding UI logic and data flow were analysed to understand how the system collects, displays, and processes colour data.

{\todo Insert short explanation of the relevant functions, e.g., where and how OBJ import populates the UI.}


Because 3MF and OBJ formats are handled through different pipelines, several adjustments were required. 
The imported colours were collected and converted into a vector<RGBA> structure. 
Since 3MF represents colours in hexadecimal form, a conversion to a floating-point representation was necessary to match the expected input format of the UI dialog (four decimal values corresponding to RGBA). 
Existing code segments for colour conversion within the \textit{OrcaSlicer} repository were reused to ensure consistency.

The dialog also required additional information about the current extruder configuration, including the colour of each extruder, the ID of the primary extruder, and associated filament identifiers. 
Once provided, the dialog correctly displayed the imported colours and allowed manual mapping to filament presets.

Although additional analysis and implementation steps had been intended, unforeseen technical challenges and limited time prevented their completion within the scope of the project.